\chapter{Unit 2: Voltage and Current Divider Laws \& Prototyping Measurements}

\section{Section A: Multiple Choice Practice Questions (60 Curated Problems)}

\noindent The following 60 multiple-choice questions reflect the official workshop evaluation bank, categorized by fundamental learning outcome and laboratory testing difficulty.\vspace{3mm}

\noindent\textbf{Q1. What is the main purpose of a voltage divider circuit?} \hfill [Tag: Validation and analysis of the Voltage Divider Law | Easy]
\begin{enumerate}[label=(\alph*),topsep=2pt,itemsep=1pt,partopsep=0pt,parsep=0pt]
    \item {To store electrical charge}
    \item {To obtain a fraction of an input voltage}
    \item {To convert DC into AC}
    \item {To increase the input current}
\end{enumerate}
\vspace{2mm}

\noindent\textbf{Q2. Which components are commonly used to form a basic voltage divider?} \hfill [Tag: Validation and analysis of the Voltage Divider Law | Easy]
\begin{enumerate}[label=(\alph*),topsep=2pt,itemsep=1pt,partopsep=0pt,parsep=0pt]
    \item {Two resistors in series}
    \item {Two diodes in series}
    \item {Two capacitors in parallel}
    \item {Two inductors in parallel}
\end{enumerate}
\vspace{2mm}

\noindent\textbf{Q3. For two series resistors, which formula gives the output voltage across $R_2$?} \hfill [Tag: Validation and analysis of the Voltage Divider Law | Easy]
\begin{enumerate}[label=(\alph*),topsep=2pt,itemsep=1pt,partopsep=0pt,parsep=0pt]
    \item {$V_{\mathrm{out}} = \frac{V_{\mathrm{in}}}{R_1 + R_2}$}
    \item {$V_{\mathrm{out}} = V_{\mathrm{in}}(R_1 + R_2)$}
    \item {$V_{\mathrm{out}} = V_{\mathrm{in}} \frac{R_1}{R_1 + R_2}$}
    \item {$V_{\mathrm{out}} = V_{\mathrm{in}} \frac{R_2}{R_1 + R_2}$}
\end{enumerate}
\vspace{2mm}

\noindent\textbf{Q4. If $R_1 = R_2$ in a voltage divider, what is the output voltage across either resistor?} \hfill [Tag: Validation and analysis of the Voltage Divider Law | Easy]
\begin{enumerate}[label=(\alph*),topsep=2pt,itemsep=1pt,partopsep=0pt,parsep=0pt]
    \item {$V_{\mathrm{in}}$}
    \item {$\frac{V_{\mathrm{in}}}{2}$}
    \item {$0$}
    \item {$2 V_{\mathrm{in}}$}
\end{enumerate}
\vspace{2mm}

\noindent\textbf{Q5. A voltage divider has V in = 12 V , R 1 = 2 k $\Omega$ , and R 2 = 4 k $\Omega$ . What is V o u t across R 2 ?} \hfill [Tag: Validation and analysis of the Voltage Divider Law | Easy]
\begin{enumerate}[label=(\alph*),topsep=2pt,itemsep=1pt,partopsep=0pt,parsep=0pt]
    \item {2 V}
    \item {8 V}
    \item {12 V}
    \item {4 V}
\end{enumerate}
\vspace{2mm}

\noindent\textbf{Q6. In a series voltage divider, which resistor has the larger voltage drop?} \hfill [Tag: Validation and analysis of the Voltage Divider Law | Easy]
\begin{enumerate}[label=(\alph*),topsep=2pt,itemsep=1pt,partopsep=0pt,parsep=0pt]
    \item {The resistor with the larger resistance}
    \item {Both always have equal drops}
    \item {The resistor connected to ground}
    \item {The resistor with the smaller resistance}
\end{enumerate}
\vspace{2mm}

\noindent\textbf{Q7. What happens to the output voltage across R 2 if R 2 increases while R 1 and V in remain constant?} \hfill [Tag: Validation and analysis of the Voltage Divider Law | Easy]
\begin{enumerate}[label=(\alph*),topsep=2pt,itemsep=1pt,partopsep=0pt,parsep=0pt]
    \item {It becomes zero}
    \item {It increases}
    \item {It decreases}
    \item {It remains unchanged}
\end{enumerate}
\vspace{2mm}

\noindent\textbf{Q8. What is the total resistance of two series resistors $R_1$ and $R_2$?} \hfill [Tag: Validation and analysis of the Voltage Divider Law | Easy]
\begin{enumerate}[label=(\alph*),topsep=2pt,itemsep=1pt,partopsep=0pt,parsep=0pt]
    \item {$R_1 R_2$}
    \item {$R_1 - R_2$}
    \item {$\frac{R_1}{R_2}$}
    \item {$R_1 + R_2$}
\end{enumerate}
\vspace{2mm}

\noindent\textbf{Q9. How can the voltage divider law be experimentally validated?} \hfill [Tag: Validation and analysis of the Voltage Divider Law | Easy]
\begin{enumerate}[label=(\alph*),topsep=2pt,itemsep=1pt,partopsep=0pt,parsep=0pt]
    \item {Measure only the resistor colors}
    \item {Disconnect the voltage source first}
    \item {Compare calculated and measured output voltages}
    \item {Replace both resistors with wires}
\end{enumerate}
\vspace{2mm}

\noindent\textbf{Q10. Where should a voltmeter be connected to measure the divider output across R 2 ?} \hfill [Tag: Validation and analysis of the Voltage Divider Law | Easy]
\begin{enumerate}[label=(\alph*),topsep=2pt,itemsep=1pt,partopsep=0pt,parsep=0pt]
    \item {Across R 2}
    \item {Across the voltage source only}
    \item {In series with R 2}
    \item {In series with the source}
\end{enumerate}
\vspace{2mm}

\noindent\textbf{Q11. What is the main purpose of a current divider circuit?} \hfill [Tag: Validation and analysis of the Current Divider Law | Easy]
\begin{enumerate}[label=(\alph*),topsep=2pt,itemsep=1pt,partopsep=0pt,parsep=0pt]
    \item {To split current among parallel branches}
    \item {To combine separate voltage sources}
    \item {To store current in a capacitor}
    \item {To increase total resistance only}
\end{enumerate}
\vspace{2mm}

\noindent\textbf{Q12. Which connection is required for a basic current divider?} \hfill [Tag: Validation and analysis of the Current Divider Law | Easy]
\begin{enumerate}[label=(\alph*),topsep=2pt,itemsep=1pt,partopsep=0pt,parsep=0pt]
    \item {Resistors connected in parallel}
    \item {Resistors connected without a source}
    \item {Resistors connected to open circuits}
    \item {Resistors connected in series}
\end{enumerate}
\vspace{2mm}

\noindent\textbf{Q13. For two parallel resistors, which formula gives the current through $R_1$?} \hfill [Tag: Validation and analysis of the Current Divider Law | Easy]
\begin{enumerate}[label=(\alph*),topsep=2pt,itemsep=1pt,partopsep=0pt,parsep=0pt]
    \item {$I_1 = I_T \frac{R_2}{R_1 + R_2}$}
    \item {$I_1 = \frac{I_T}{R_1 + R_2}$}
    \item {$I_1 = I_T(R_1 + R_2)$}
    \item {$I_1 = I_T \frac{R_1}{R_1 + R_2}$}
\end{enumerate}
\vspace{2mm}

\noindent\textbf{Q14. In a parallel circuit, which branch carries more current?} \hfill [Tag: Validation and analysis of the Current Divider Law | Easy]
\begin{enumerate}[label=(\alph*),topsep=2pt,itemsep=1pt,partopsep=0pt,parsep=0pt]
    \item {Both branches always carry zero current}
    \item {The branch with larger resistance}
    \item {The branch with no voltage}
    \item {The branch with smaller resistance}
\end{enumerate}
\vspace{2mm}

\noindent\textbf{Q15. If two parallel resistors have equal resistance, how is the total current divided?} \hfill [Tag: Validation and analysis of the Current Divider Law | Easy]
\begin{enumerate}[label=(\alph*),topsep=2pt,itemsep=1pt,partopsep=0pt,parsep=0pt]
    \item {All current flows through one branch}
    \item {The current doubles in both branches}
    \item {The current divides equally}
    \item {No current enters either branch}
\end{enumerate}
\vspace{2mm}

\noindent\textbf{Q16. A total current of 6 A enters two equal parallel resistors. What current flows through each resistor?} \hfill [Tag: Validation and analysis of the Current Divider Law | Easy]
\begin{enumerate}[label=(\alph*),topsep=2pt,itemsep=1pt,partopsep=0pt,parsep=0pt]
    \item {6 A}
    \item {2 A}
    \item {1 A}
    \item {3 A}
\end{enumerate}
\vspace{2mm}

\noindent\textbf{Q17. A current divider has I T = 6 A , R 1 = 2 $\Omega$ , and R 2 = 4 $\Omega$ . What is the current through R 1 ?} \hfill [Tag: Validation and analysis of the Current Divider Law | Easy]
\begin{enumerate}[label=(\alph*),topsep=2pt,itemsep=1pt,partopsep=0pt,parsep=0pt]
    \item {2 A}
    \item {6 A}
    \item {1 A}
    \item {4 A}
\end{enumerate}
\vspace{2mm}

\noindent\textbf{Q18. What is the relationship between total current and branch currents in a two-branch parallel circuit?} \hfill [Tag: Validation and analysis of the Current Divider Law | Easy]
\begin{enumerate}[label=(\alph*),topsep=2pt,itemsep=1pt,partopsep=0pt,parsep=0pt]
    \item {$I_T = I_1 + I_2$}
    \item {$I_T = I_1 - I_2$}
    \item {$I_T = \frac{I_1}{I_2}$}
    \item {$I_T = I_1 I_2$}
\end{enumerate}
\vspace{2mm}

\noindent\textbf{Q19. How can the current divider law be experimentally validated?} \hfill [Tag: Validation and analysis of the Current Divider Law | Easy]
\begin{enumerate}[label=(\alph*),topsep=2pt,itemsep=1pt,partopsep=0pt,parsep=0pt]
    \item {Measure only the supply voltage}
    \item {Connect the branches in series}
    \item {Compare calculated and measured branch currents}
    \item {Remove the parallel resistors}
\end{enumerate}
\vspace{2mm}

\noindent\textbf{Q20. Where should an ammeter be connected to measure the current through one branch?} \hfill [Tag: Validation and analysis of the Current Divider Law | Easy]
\begin{enumerate}[label=(\alph*),topsep=2pt,itemsep=1pt,partopsep=0pt,parsep=0pt]
    \item {Across the voltage source}
    \item {Across the entire circuit}
    \item {In series with the branch}
    \item {In parallel with the branch}
\end{enumerate}
\vspace{2mm}

\noindent\textbf{Q21. A 12 V source is connected across two series resistors, R 1 = 2 k $\Omega$ and R 2 = 4 k $\Omega$ . What is the voltage across R 2 ?} \hfill [Tag: Validation and analysis of the Voltage Divider Law | Medium]
\begin{enumerate}[label=(\alph*),topsep=2pt,itemsep=1pt,partopsep=0pt,parsep=0pt]
    \item {6 V}
    \item {4 V}
    \item {10 V}
    \item {8 V}
\end{enumerate}
\vspace{2mm}

\noindent\textbf{Q22. A divider uses R 1 = 3 k $\Omega$ and R 2 = 2 k $\Omega$ across a 10 V supply. If the measured output across R 2 is 3.8 V , what is the percentage error relative to the theoretical value?} \hfill [Tag: Validation and analysis of the Voltage Divider Law | Medium]
\begin{enumerate}[label=(\alph*),topsep=2pt,itemsep=1pt,partopsep=0pt,parsep=0pt]
    \item {7.5\%}
    \item {2.5\%}
    \item {5.0\%}
    \item {10.0\%}
\end{enumerate}
\vspace{2mm}

\noindent\textbf{Q23. A 12 V divider has R 1 = 2 k $\Omega$ and R 2 = 6 k $\Omega$ . A 6 k $\Omega$ load is connected in parallel with R 2 . What is the loaded output voltage?} \hfill [Tag: Validation and analysis of the Voltage Divider Law | Medium]
\begin{enumerate}[label=(\alph*),topsep=2pt,itemsep=1pt,partopsep=0pt,parsep=0pt]
    \item {8.0 V}
    \item {6.0 V}
    \item {7.2 V}
    \item {9.0 V}
\end{enumerate}
\vspace{2mm}

\noindent\textbf{Q24. A voltage divider must produce 5 V from a 15 V source, with the output taken across R 2 . Which resistance relationship is required?} \hfill [Tag: Validation and analysis of the Voltage Divider Law | Medium]
\begin{enumerate}[label=(\alph*),topsep=2pt,itemsep=1pt,partopsep=0pt,parsep=0pt]
    \item {R 1 = 3 R 2}
    \item {R 2 = 2 R 1}
    \item {R 1 = R 2}
    \item {R 1 = 2 R 2}
\end{enumerate}
\vspace{2mm}

\noindent\textbf{Q25. Two series resistors, 1 k $\Omega$ and 2 k $\Omega$ , are connected to a 9 V source. Which pair of measured voltage drops would validate the voltage divider law?} \hfill [Tag: Validation and analysis of the Voltage Divider Law | Medium]
\begin{enumerate}[label=(\alph*),topsep=2pt,itemsep=1pt,partopsep=0pt,parsep=0pt]
    \item {6 V and 3 V}
    \item {2 V and 7 V}
    \item {4 V and 5 V}
    \item {3 V and 6 V}
\end{enumerate}
\vspace{2mm}

\noindent\textbf{Q26. Which voltmeter characteristic most effectively reduces loading error when experimentally validating a voltage divider?} \hfill [Tag: Validation and analysis of the Voltage Divider Law | Medium]
\begin{enumerate}[label=(\alph*),topsep=2pt,itemsep=1pt,partopsep=0pt,parsep=0pt]
    \item {Low input resistance}
    \item {Low voltage range}
    \item {High current range}
    \item {High input resistance}
\end{enumerate}
\vspace{2mm}

\noindent\textbf{Q27. A 10 V divider uses two nominal 10 k $\Omega$ resistors with $\pm$ 1\% tolerance. What is the approximate minimum output across the lower resistor due only to resistor tolerance?} \hfill [Tag: Validation and analysis of the Voltage Divider Law | Medium]
\begin{enumerate}[label=(\alph*),topsep=2pt,itemsep=1pt,partopsep=0pt,parsep=0pt]
    \item {4.90 V}
    \item {5.00 V}
    \item {4.95 V}
    \item {5.05 V}
\end{enumerate}
\vspace{2mm}

\noindent\textbf{Q28. An 18 V divider has R 1 = 3 k $\Omega$ and R 2 = 6 k $\Omega$ . A 4 k $\Omega$ load is connected to its output. Using the divider's Thevenin equivalent, what is the loaded output voltage?} \hfill [Tag: Validation and analysis of the Voltage Divider Law | Medium]
\begin{enumerate}[label=(\alph*),topsep=2pt,itemsep=1pt,partopsep=0pt,parsep=0pt]
    \item {12 V}
    \item {6 V}
    \item {10 V}
    \item {8 V}
\end{enumerate}
\vspace{2mm}

\noindent\textbf{Q29. A 20 V source supplies series resistors R 1 = 2 k $\Omega$ and R 2 = 3 k $\Omega$ . What power is dissipated by R 2 ?} \hfill [Tag: Validation and analysis of the Voltage Divider Law | Medium]
\begin{enumerate}[label=(\alph*),topsep=2pt,itemsep=1pt,partopsep=0pt,parsep=0pt]
    \item {32 mW}
    \item {60 mW}
    \item {40 mW}
    \item {48 mW}
\end{enumerate}
\vspace{2mm}

\noindent\textbf{Q30. The theoretical output of a divider is 6.00 V , while the measured output is 5.82 V . What is the magnitude of the percentage deviation?} \hfill [Tag: Validation and analysis of the Voltage Divider Law | Medium]
\begin{enumerate}[label=(\alph*),topsep=2pt,itemsep=1pt,partopsep=0pt,parsep=0pt]
    \item {3\%}
    \item {4\%}
    \item {5\%}
    \item {2\%}
\end{enumerate}
\vspace{2mm}

\noindent\textbf{Q31. A total current of 6 A enters two parallel resistors of 6 $\Omega$ and 3 $\Omega$ . What current flows through the 6 $\Omega$ resistor?} \hfill [Tag: Validation and analysis of the Current Divider Law | Medium]
\begin{enumerate}[label=(\alph*),topsep=2pt,itemsep=1pt,partopsep=0pt,parsep=0pt]
    \item {4 A}
    \item {1 A}
    \item {2 A}
    \item {3 A}
\end{enumerate}
\vspace{2mm}

\noindent\textbf{Q32. Two parallel branches have conductances of 2 mS and 6 mS . If the total current is 16 mA , what current flows through the 6 mS branch?} \hfill [Tag: Validation and analysis of the Current Divider Law | Medium]
\begin{enumerate}[label=(\alph*),topsep=2pt,itemsep=1pt,partopsep=0pt,parsep=0pt]
    \item {12 mA}
    \item {8 mA}
    \item {4 mA}
    \item {14 mA}
\end{enumerate}
\vspace{2mm}

\noindent\textbf{Q33. A total current of 6 mA enters three parallel resistors of 2 k $\Omega$ , 3 k $\Omega$ , and 6 k $\Omega$ . What current flows through the 3 k $\Omega$ branch?} \hfill [Tag: Validation and analysis of the Current Divider Law | Medium]
\begin{enumerate}[label=(\alph*),topsep=2pt,itemsep=1pt,partopsep=0pt,parsep=0pt]
    \item {2 mA}
    \item {3 mA}
    \item {4 mA}
    \item {1 mA}
\end{enumerate}
\vspace{2mm}

\noindent\textbf{Q34. A current of 10 mA enters a parallel combination of 1 k $\Omega$ and 4 k $\Omega$ . What is the voltage across the parallel network?} \hfill [Tag: Validation and analysis of the Current Divider Law | Medium]
\begin{enumerate}[label=(\alph*),topsep=2pt,itemsep=1pt,partopsep=0pt,parsep=0pt]
    \item {10 V}
    \item {4 V}
    \item {6 V}
    \item {8 V}
\end{enumerate}
\vspace{2mm}

\noindent\textbf{Q35. During validation of a parallel circuit, the total current is measured as 8.0 mA , while the branch currents are 5.1 mA and 2.8 mA . What is the percentage discrepancy relative to the total current?} \hfill [Tag: Validation and analysis of the Current Divider Law | Medium]
\begin{enumerate}[label=(\alph*),topsep=2pt,itemsep=1pt,partopsep=0pt,parsep=0pt]
    \item {1.25\%}
    \item {0.80\%}
    \item {2.50\%}
    \item {3.75\%}
\end{enumerate}
\vspace{2mm}

\noindent\textbf{Q36. A total current of 12 mA divides between parallel resistors of 2 k $\Omega$ and 4 k $\Omega$ . What power is dissipated by the 4 k $\Omega$ resistor?} \hfill [Tag: Validation and analysis of the Current Divider Law | Medium]
\begin{enumerate}[label=(\alph*),topsep=2pt,itemsep=1pt,partopsep=0pt,parsep=0pt]
    \item {96 mW}
    \item {48 mW}
    \item {64 mW}
    \item {32 mW}
\end{enumerate}
\vspace{2mm}

\noindent\textbf{Q37. A total current of 5 mA enters two parallel branches. The first branch contains 2 k $\Omega$ and carries 3 mA . What is the resistance of the second branch?} \hfill [Tag: Validation and analysis of the Current Divider Law | Medium]
\begin{enumerate}[label=(\alph*),topsep=2pt,itemsep=1pt,partopsep=0pt,parsep=0pt]
    \item {1 k $\Omega$}
    \item {2 k $\Omega$}
    \item {4 k $\Omega$}
    \item {3 k $\Omega$}
\end{enumerate}
\vspace{2mm}

\noindent\textbf{Q38. Under which circuit condition can the current divider law be directly applied?} \hfill [Tag: Validation and analysis of the Current Divider Law | Medium]
\begin{enumerate}[label=(\alph*),topsep=2pt,itemsep=1pt,partopsep=0pt,parsep=0pt]
    \item {Resistors dissipate the same power}
    \item {Resistors carry the same series current}
    \item {Resistors share the same two nodes}
    \item {Resistors have equal voltage ratings}
\end{enumerate}
\vspace{2mm}

\noindent\textbf{Q39. An ideal voltage source supplies parallel branches of 2 k $\Omega$ and 3 k $\Omega$ . An ammeter with 100 $\Omega$ internal resistance is inserted into the 2 k $\Omega$ branch. What happens to the measured circuit currents?} \hfill [Tag: Validation and analysis of the Current Divider Law | Medium]
\begin{enumerate}[label=(\alph*),topsep=2pt,itemsep=1pt,partopsep=0pt,parsep=0pt]
    \item {Both branch currents become equal}
    \item {The second and total currents increase}
    \item {Both branch currents remain unchanged}
    \item {The first and total currents decrease}
\end{enumerate}
\vspace{2mm}

\noindent\textbf{Q40. A 12 mA current divides between a 4 k $\Omega$ resistor and an unknown parallel resistor. If the 4 k $\Omega$ branch carries 3 mA , what is the unknown resistance?} \hfill [Tag: Validation and analysis of the Current Divider Law | Medium]
\begin{enumerate}[label=(\alph*),topsep=2pt,itemsep=1pt,partopsep=0pt,parsep=0pt]
    \item {2.00 k $\Omega$}
    \item {3.00 k $\Omega$}
    \item {1.00 k $\Omega$}
    \item {1.33 k $\Omega$}
\end{enumerate}
\vspace{2mm}

\noindent\textbf{Q41. An 18 V source feeds R 1 = 2 k $\Omega$ in series with R 2 = 4 k $\Omega$ . A 12 k $\Omega$ load is connected across R 2 . What is the loaded output voltage across the load?} \hfill [Tag: Validation and analysis of the Voltage Divider Law | Hard]
\begin{enumerate}[label=(\alph*),topsep=2pt,itemsep=1pt,partopsep=0pt,parsep=0pt]
    \item {6.00 V}
    \item {7.20 V}
    \item {8.00 V}
    \item {9.00 V}
\end{enumerate}
\vspace{2mm}

\noindent\textbf{Q42. A 10 V source is connected to R 1 = 100 k $\Omega$ and R 2 = 900 k $\Omega$ . A voltmeter with input resistance 1 M $\Omega$ measures the voltage across R 2 . What reading should be expected?} \hfill [Tag: Validation and analysis of the Voltage Divider Law | Hard]
\begin{enumerate}[label=(\alph*),topsep=2pt,itemsep=1pt,partopsep=0pt,parsep=0pt]
    \item {7.50 V}
    \item {8.75 V}
    \item {8.26 V}
    \item {9.00 V}
\end{enumerate}
\vspace{2mm}

\noindent\textbf{Q43. For a divider with fixed input voltage and fixed total resistance R 1 + R 2 = 10 k $\Omega$ , which choice maximizes the power dissipated in R 2 ?} \hfill [Tag: Validation and analysis of the Voltage Divider Law | Hard]
\begin{enumerate}[label=(\alph*),topsep=2pt,itemsep=1pt,partopsep=0pt,parsep=0pt]
    \item {R 1 = 5 k $\Omega$ , R 2 = 5 k $\Omega$}
    \item {R 1 = 4 k $\Omega$ , R 2 = 6 k $\Omega$}
    \item {R 1 = 2 k $\Omega$ , R 2 = 8 k $\Omega$}
    \item {R 1 = 8 k $\Omega$ , R 2 = 2 k $\Omega$}
\end{enumerate}
\vspace{2mm}

\noindent\textbf{Q44. A divider has R 1 = 1 k $\Omega$ and R 2 = 2 k $\Omega$ , each with a tolerance of $\pm$ 5\% . What are the worst-case limits of V o / V i when the output is taken across R 2 ?} \hfill [Tag: Validation and analysis of the Voltage Divider Law | Hard]
\begin{enumerate}[label=(\alph*),topsep=2pt,itemsep=1pt,partopsep=0pt,parsep=0pt]
    \item {0.633 to 0.667}
    \item {0.615 to 0.680}
    \item {0.600 to 0.700}
    \item {0.603 to 0.689}
\end{enumerate}
\vspace{2mm}

\noindent\textbf{Q45. A 12 V source feeds R 1 = 3 k $\Omega$ and R 2 = 3 k $\Omega$ . A 6 k $\Omega$ load is connected across R 2 . What is the output voltage?} \hfill [Tag: Validation and analysis of the Voltage Divider Law | Hard]
\begin{enumerate}[label=(\alph*),topsep=2pt,itemsep=1pt,partopsep=0pt,parsep=0pt]
    \item {4.00 V}
    \item {3.60 V}
    \item {6.00 V}
    \item {4.80 V}
\end{enumerate}
\vspace{2mm}

\noindent\textbf{Q46. A divider is driven by an ideal 5 V source through a source resistance of 50 $\Omega$ , followed by R 1 = 950 $\Omega$ and R 2 = 1 k $\Omega$ . What voltage appears across R 2 ?} \hfill [Tag: Validation and analysis of the Voltage Divider Law | Hard]
\begin{enumerate}[label=(\alph*),topsep=2pt,itemsep=1pt,partopsep=0pt,parsep=0pt]
    \item {2.63 V}
    \item {2.44 V}
    \item {2.56 V}
    \item {2.50 V}
\end{enumerate}
\vspace{2mm}

\noindent\textbf{Q47. A divider uses R 1 = 2 k $\Omega$ , R 2 = 3 k $\Omega$ , and a load R L . With V i = 15 V , the desired loaded output is 6 V . What value of R L is required?} \hfill [Tag: Validation and analysis of the Voltage Divider Law | Hard]
\begin{enumerate}[label=(\alph*),topsep=2pt,itemsep=1pt,partopsep=0pt,parsep=0pt]
    \item {2.40 k $\Omega$}
    \item {1.20 k $\Omega$}
    \item {1.80 k $\Omega$}
    \item {3.60 k $\Omega$}
\end{enumerate}
\vspace{2mm}

\noindent\textbf{Q48. A voltage divider has R 1 = 4 k $\Omega$ and R 2 = 6 k $\Omega$ . If a finite-input-resistance voltmeter changes the measured output from the ideal value by less than 1\% , what is the minimum approximate meter resistance?} \hfill [Tag: Validation and analysis of the Voltage Divider Law | Hard]
\begin{enumerate}[label=(\alph*),topsep=2pt,itemsep=1pt,partopsep=0pt,parsep=0pt]
    \item {60 k $\Omega$}
    \item {600 k $\Omega$}
    \item {300 k $\Omega$}
    \item {120 k $\Omega$}
\end{enumerate}
\vspace{2mm}

\noindent\textbf{Q49. For an unloaded divider, V o = V i R 2 / ( R 1 + R 2 ) . If R 1 is increased by 10\% while R 2 remains fixed, with R 1 = R 2 initially, what is the percentage change in V o ?} \hfill [Tag: Validation and analysis of the Voltage Divider Law | Hard]
\begin{enumerate}[label=(\alph*),topsep=2pt,itemsep=1pt,partopsep=0pt,parsep=0pt]
    \item {A decrease of approximately 9.1\%}
    \item {An increase of approximately 9.1\%}
    \item {A decrease of approximately 4.8\%}
    \item {An increase of approximately 4.8\%}
\end{enumerate}
\vspace{2mm}

\noindent\textbf{Q50. A practical divider is connected to a 20 V source with internal resistance 100 $\Omega$ . The divider consists of R 1 = 900 $\Omega$ and R 2 = 1 k $\Omega$ . What is the actual output across R 2 ?} \hfill [Tag: Validation and analysis of the Voltage Divider Law | Hard]
\begin{enumerate}[label=(\alph*),topsep=2pt,itemsep=1pt,partopsep=0pt,parsep=0pt]
    \item {9.52 V}
    \item {9.00 V}
    \item {10.00 V}
    \item {10.53 V}
\end{enumerate}
\vspace{2mm}

\noindent\textbf{Q51. A total current of 6 mA enters three parallel resistors of 2 k $\Omega$ , 3 k $\Omega$ , and 6 k $\Omega$ . What current flows through the 2 k $\Omega$ branch?} \hfill [Tag: Validation and analysis of the Current Divider Law | Hard]
\begin{enumerate}[label=(\alph*),topsep=2pt,itemsep=1pt,partopsep=0pt,parsep=0pt]
    \item {1.0 mA}
    \item {2.0 mA}
    \item {4.0 mA}
    \item {3.0 mA}
\end{enumerate}
\vspace{2mm}

\noindent\textbf{Q52. A 10 mA current source feeds parallel branches of 4 k $\Omega$ and 6 k $\Omega$ . What are the currents in the 4 k $\Omega$ and 6 k $\Omega$ branches, respectively?} \hfill [Tag: Validation and analysis of the Current Divider Law | Hard]
\begin{enumerate}[label=(\alph*),topsep=2pt,itemsep=1pt,partopsep=0pt,parsep=0pt]
    \item {7.5 mA and 2.5 mA}
    \item {5 mA and 5 mA}
    \item {6 mA and 4 mA}
    \item {4 mA and 6 mA}
\end{enumerate}
\vspace{2mm}

\noindent\textbf{Q53. A 3 mA ideal current source feeds three parallel resistances: 10 k $\Omega$ , 20 k $\Omega$ , and 30 k $\Omega$ . What current flows through the 10 k $\Omega$ branch?} \hfill [Tag: Validation and analysis of the Current Divider Law | Hard]
\begin{enumerate}[label=(\alph*),topsep=2pt,itemsep=1pt,partopsep=0pt,parsep=0pt]
    \item {2.00 mA}
    \item {1.64 mA}
    \item {1.80 mA}
    \item {1.09 mA}
\end{enumerate}
\vspace{2mm}

\noindent\textbf{Q54. A 12 A current enters two parallel paths. One path is 4 $\Omega$ ; the other consists of 6 $\Omega$ in series with 2 $\Omega$ . What current flows through the 2 $\Omega$ resistor?} \hfill [Tag: Validation and analysis of the Current Divider Law | Hard]
\begin{enumerate}[label=(\alph*),topsep=2pt,itemsep=1pt,partopsep=0pt,parsep=0pt]
    \item {8 A}
    \item {4 A}
    \item {6 A}
    \item {3 A}
\end{enumerate}
\vspace{2mm}

\noindent\textbf{Q55. Two parallel branches have resistances 8 $\Omega$ and 12 $\Omega$ . If the current in the 8 $\Omega$ branch is 3 A , what is the total source current?} \hfill [Tag: Validation and analysis of the Current Divider Law | Hard]
\begin{enumerate}[label=(\alph*),topsep=2pt,itemsep=1pt,partopsep=0pt,parsep=0pt]
    \item {4.5 A}
    \item {5.0 A}
    \item {7.5 A}
    \item {4.0 A}
\end{enumerate}
\vspace{2mm}

\noindent\textbf{Q56. A current divider has two branches, R 1 = 5 k $\Omega$ and R 2 = 15 k $\Omega$ . If the total current is 8 mA , what is the current through R 2 ?} \hfill [Tag: Validation and analysis of the Current Divider Law | Hard]
\begin{enumerate}[label=(\alph*),topsep=2pt,itemsep=1pt,partopsep=0pt,parsep=0pt]
    \item {6.0 mA}
    \item {1.0 mA}
    \item {4.0 mA}
    \item {2.0 mA}
\end{enumerate}
\vspace{2mm}

\noindent\textbf{Q57. A 2 A source supplies two parallel branches. Branch A has 3 $\Omega$ and branch B has 6 $\Omega$ . An ammeter with 0.5 $\Omega$ burden resistance is inserted only in branch A. What current does the ammeter indicate?} \hfill [Tag: Validation and analysis of the Current Divider Law | Hard]
\begin{enumerate}[label=(\alph*),topsep=2pt,itemsep=1pt,partopsep=0pt,parsep=0pt]
    \item {0.80 A}
    \item {1.00 A}
    \item {1.33 A}
    \item {1.14 A}
\end{enumerate}
\vspace{2mm}

\noindent\textbf{Q58. Three equal resistors are connected in parallel to an ideal current source. One resistor is then short-circuited by an ideal wire. What happens to the current in each resistor?} \hfill [Tag: Validation and analysis of the Current Divider Law | Hard]
\begin{enumerate}[label=(\alph*),topsep=2pt,itemsep=1pt,partopsep=0pt,parsep=0pt]
    \item {All three retain one-third of the source current}
    \item {The current divides equally between the two unshorted resistors}
    \item {The shorted resistor carries half the current; the others share half}
    \item {The shorted resistor carries all current; the others carry zero}
\end{enumerate}
\vspace{2mm}

\noindent\textbf{Q59. A total current $I_T$ enters two parallel resistors $R_1$ and $R_2$. Which expression gives the current through $R_1$ and remains valid when $R_1 \neq R_2$?} \hfill [Tag: Validation and analysis of the Current Divider Law | Hard]
\begin{enumerate}[label=(\alph*),topsep=2pt,itemsep=1pt,partopsep=0pt,parsep=0pt]
    \item {$I_1 = I_T \frac{R_1 R_2}{R_1 + R_2}$}
    \item {$I_1 = I_T \frac{R_1 + R_2}{R_2}$}
    \item {$I_1 = I_T \frac{R_1}{R_1 + R_2}$}
    \item {$I_1 = I_T \frac{R_2}{R_1 + R_2}$}
\end{enumerate}
\vspace{2mm}

\noindent\textbf{Q60. A 6 mA current source feeds two parallel branches: R 1 = 2 k $\Omega$ and R 2 = 4 k $\Omega$ . A 6 k $\Omega$ resistor is connected in parallel with R 2 . What current flows through the original R 1 branch?} \hfill [Tag: Validation and analysis of the Current Divider Law | Hard]
\begin{enumerate}[label=(\alph*),topsep=2pt,itemsep=1pt,partopsep=0pt,parsep=0pt]
    \item {3.00 mA}
    \item {4.00 mA}
    \item {3.60 mA}
    \item {2.00 mA}
\end{enumerate}
\vspace{2mm}


\newpage
\section{Section B: Comprehensive Subjective Problem with Detailed Analytical Rubric}

\begin{subjectiveq}{2}{1}
A series circuit consists of $V_{\mathrm{in}} = 12\,\mathrm{V}$, $R_1 = 10\,\mathrm{k}\Omega$, and $R_2 = 10\,\mathrm{k}\Omega$. (a) Calculate the theoretical unloaded voltage across $R_2$. (b) If an analog voltmeter with internal resistance $R_m = 20\,\mathrm{k}\Omega$ is connected across $R_2$, determine the measured voltage and the resulting percentage loading error. (c) What is the required meter impedance to keep loading error below $1\%$?
\end{subjectiveq}

\begin{solbox}
\textbf{(a) Unloaded Output Voltage:}\\
By the Voltage Divider Law:\\
\[ V_2 = V_{\mathrm{in}} \frac{R_2}{R_1 + R_2} = 12\,\mathrm{V} \times \frac{10\,\mathrm{k}\Omega}{10\,\mathrm{k}\Omega + 10\,\mathrm{k}\Omega} = 12 \times \frac{1}{2} = 6.00\,\mathrm{V} \]
\textbf{(b) Loaded Output Voltage:}\\
Connecting the meter across $R_2$ places $R_m$ in parallel with $R_2$:\\
\[ R_{2,\mathrm{eq}} = R_2 \parallel R_m = \frac{R_2 \cdot R_m}{R_2 + R_m} = \frac{10 \times 20}{10 + 20}\,\mathrm{k}\Omega = \frac{200}{30} = 6.67\,\mathrm{k}\Omega \]
The loaded divider voltage is:\\
\[ V_{2,\mathrm{meas}} = V_{\mathrm{in}} \frac{R_{2,\mathrm{eq}}}{R_1 + R_{2,\mathrm{eq}}} = 12 \times \frac{6.67}{10 + 6.67} = 12 \times \frac{6.67}{16.67} = 4.80\,\mathrm{V} \]
Loading Error:\\
\[ \%\text{ Error} = \frac{V_{\mathrm{true}} - V_{\mathrm{meas}}}{V_{\mathrm{true}}} \times 100\% = \frac{6.00 - 4.80}{6.00} \times 100\% = 20.0\% \]
\textbf{(c) High-Impedance Condition for $\le 1\%$ Error:}\\
For error $\le 1\%$, $V_{\mathrm{meas}} \ge 5.94\,\mathrm{V}$. Setting $V_{\mathrm{meas}} = 12 \frac{R_{2,\mathrm{eq}}}{10 + R_{2,\mathrm{eq}}} \ge 5.94$ yields $R_{2,\mathrm{eq}} \ge 9.80\,\mathrm{k}\Omega$, requiring $R_m \ge \frac{10 \times 9.80}{10 - 9.80} \approx 490\,\mathrm{k}\Omega$. Modern DMMs have $R_m = 10\,\mathrm{M}\Omega$, ensuring negligible error ($< 0.05\%$).
\end{solbox}

\begin{examinertip}
Examiner Rubric (10 Marks):\\
- Unloaded voltage derivation \& correct value: 2 Marks\\
- Equivalent parallel resistance $R_2 \parallel R_m$: 3 Marks\\
- Loaded voltage computation \& percentage error: 3 Marks\\
- Analysis of meter impedance threshold for precision measurement: 2 Marks
\end{examinertip}


\section{Section C: Laboratory Viva-Voce Questions with Authoritative Solutions}

\noindent\textbf{Viva Question:} State the dual relationship between the Voltage Divider Law and the Current Divider Law.\\
\textbf{Technical Answer:} The Voltage Divider Law applies to series circuits where current is constant and source voltage distributes proportionally to individual resistances: $V_x = V_T (R_x / R_{\mathrm{total}})$. The Current Divider Law applies to parallel circuits where voltage is constant and total current distributes proportionally to individual conductances: $I_x = I_T (G_x / G_{\mathrm{total}}) = I_T (R_p / R_x)$.
\vspace{3mm}

\noindent\textbf{Viva Question:} Why does a potentiometer act as a variable voltage divider?\\
\textbf{Technical Answer:} A potentiometer consists of a resistive track of fixed length with a sliding wiper contact. The wiper divides the track into two series resistors $R_1$ and $R_2$ whose sum is constant ($R_1 + R_2 = R_{\mathrm{pot}}$). Moving the wiper adjusts the ratio $R_2 / (R_1 + R_2)$ continuously from 0 to 1, producing an output voltage from $0\text{ V}$ to $V_{CC}$.
\vspace{3mm}

\noindent\textbf{Viva Question:} What is ammeter burden voltage?\\
\textbf{Technical Answer:} Ammeter burden voltage is the small voltage drop developed across the internal current-sensing shunt resistor when current flows through the meter. In low-voltage circuits, this burden voltage reduces total branch current, introducing measurement error.
\vspace{3mm}


\section{Section D: Workshop Troubleshooting \& Fault Diagnostic Case Study}

\textbf{Troubleshooting Scenario: Voltage Divider Output Drops When Load Connected}\\
A student constructs a $5\text{V}$ to $2.5\text{V}$ reference voltage divider using two $100\,\mathrm{k}\Omega$ resistors from a $5\text{V}$ supply. When an external circuit with an input resistance of $20\,\mathrm{k}\Omega$ is attached to the output, the voltage collapses to $0.83\text{V}$.\\\\
\textbf{Diagnosis:} The student failed to account for source output impedance (Thevenin equivalent resistance $R_{th} = 100\,\mathrm{k}\Omega \parallel 100\,\mathrm{k}\Omega = 50\,\mathrm{k}\Omega$). Connecting a low load impedance ($20\,\mathrm{k}\Omega$) severely loads the divider.\\
\textbf{Remedy:} Either decrease divider resistors to $R_1 = R_2 = 200\,\Omega$ (at the cost of higher quiescent power consumption) or place an op-amp unity-gain voltage follower (buffer) between the divider and load.
\vspace{4mm}
